\begin{center}
  \textbf{ОГЛЯД ПРОБЛЕМИ}
\end{center}
Моделювання процесів теплообміну у тканинах людського тіла можна описати кількома моделями. Зокрема рівняння Пеннеса
\cite{thermal-process-tissue-pennes}. Це рівняння не моделює судини безпосередньо, а враховує їхній сукупний вплив.
Задачу теплообміну у людських тканинах у заданій області $\Omega$ можна сформулювати у вигляді:

\begin{equation}
  \label{eqn:theory_thermo_pennes_eqn}
  \nabla \cdot \lambda_\theta \nabla \theta(x) + \omega \rho_b C_b (\theta_c - \theta(x)) + q_m  = 0, x \in \Omega
\end{equation}

\noindent де $\Gamma = \partial \Omega$ -- межа області $\Omega$, $\lambda_\theta$ -- коефіцієнт теплопровідності,
$\theta$ -- шукана температура в області $\Omega$,
$\rho_b$ -- густина крові, $\omega$ -- перфузія крові через орган, $C_b$ -- питома теплоємність крові,
$\theta_c$ -- внутрішня температура тіла, $q_m$ -- метаболічне тепло у тканині. На межі задамо умови:

\begin{equation}
  \label{eqn:theory_thermo_cond_dirichlet}
  \theta(x) = \theta_{\Gamma_D}(x), x \in \Gamma_D
\end{equation}

\begin{equation}
  \label{eqn:theory_thermo_cond_neumann}
  q(x) = q_{\Gamma_N}(x), x \in \Gamma_N
\end{equation}

\begin{equation}
  \label{eqn:theory_thermo_cond_robin}
  -\lambda_\theta q(x) = h(\theta(x) - \theta_e), x \in \Gamma_R,
\end{equation}

\noindent де $q(x) = \frac{\partial \theta(x)}{\partial n}$ -- тепловий потік, $n_i$ -- компоненти вектора зовнішньої
нормалі до межі $\Gamma$, $\theta_e$ -- температура середовища, $h$ -- коефіцієнт тепловіддачі. Крайові умови
(\ref{eqn:theory_thermo_cond_dirichlet}), (\ref{eqn:theory_thermo_cond_neumann}) і (\ref{eqn:theory_thermo_cond_robin})
-- це умови Діріхле, Неймана і Робіна відповідно. Умова Робіна описує тепловий обмін тіла із навколишнім середовищем.

\begin{center}
  \textbf{ПРОПОНОВАНИЙ ПІДХІД}
\end{center}
Для розв'язування задач типу (\ref{eqn:theory_thermo_pennes_eqn}) (\ref{eqn:theory_thermo_cond_dirichlet})
(\ref{eqn:theory_thermo_cond_neumann}) (\ref{eqn:theory_thermo_cond_robin}) використаємо метод приграничних елементів
(МПГЕ)\cite{grytsko-cobem}, який належить до сімейства методів граничних елементів (МГЕ)\cite{brebbia-bem}. МПГЕ володіє
додатковим ступенем вільності, а саме шириною приграничної зони. Він дозволяє покращувати точність стаціонарних задач.

\begin{center}
  \textbf{ДИСКРЕТИЗАЦІЯ}
\end{center}
Визначимо область $G$, яка містить область $\Omega$, причому $\partial G \cap \partial \Omega = \varnothing$.  


\begin{center}
  \textbf{РОЗВ’ЯЗОК}
\end{center}

% \textbf{Актуальність теми.} Традиційне хірургічне втручання -- відкрита хірургія характеризується відносно великим
% розтином м'яких тканин. Для доступу до органу над яким хірург має провести маніпуляцію, також треба зупинити кровотечу,
% спричинену розтином. Ці дії проводять в умовах стерильності, що серйозно обмежує кількість рухів, які виконує лікар, і
% зменшує комфорт проведення операційного втручання. З іншого боку великий розтин дозволяє отримати широкий кут огляду
% місця проведення хірургічної операції. Проте така зручність супроводжується суттєвими недоліками:
% \begin{itemize}
%   \item збільшення ймовірності інфекції;
%   \item сильні больові відчуття у пацієнта\textbackslashпацієнтки після операції;
%   \item довший період реабілітації після хірургічного втручання;
%   \item великий шрам у місці розтину.
% \end{itemize}

% Головним фактором вищезгаданих недоліків є великий розтин у м'яких тканинах. Тому були створені нові методи
% малоінвазивної хірургії. Суттєва різниця полягає у розмірі отвору, зокрема у лапароскопічній операції, яка є одним із
% підвидів малоінвазивної хірургії, через отвір утворений у результаті проколу м'яких тканин, вводять інструменти. Серед
% них є ендоскоп, який виводить зображення на монітор і освітлює місце проведення операції.

% % \begin{figure}[ht!]
% %   \centering
% %   \copyrightbox[b]
% %     {\includegraphics[width=0.4\linewidth]{images/intro_laparoscopy.png}}
% %     {Medical gallery of Blausen Medical 2014}
% %   \caption{Лапароскопія}
% %   \label{fig:intro_laparoscopy}
% % \end{figure}

% Лапароскопічні операції рис. \ref{fig:intro_laparoscopy} вирішують недоліки відкритих операцій і є кроком вперед проте
% мають низку недоліків:
% \begin{itemize}
%   \item неінтуїтивне керування інструментами;
%   \item ризик ускладнень у зв'язку із обмеженим простором для операції і вищою ймовірністю пошкодити інші органи;
%   \item обмежене тактильне відчуття через використання інструментів;
%   \item довший час операції;
%   \item хірург повинен\textbackslashповинна бути у стерильній зоні протягом операції.
% \end{itemize}

% Малоінвазивна роботохірургія є наступним удосконаленням процесу проведення хірургічного втручання. Завдяки поєднанню
% сильних сторін людини і комп'ютера вдалося усунути майже всі недоліки лапароскопічних операцій. Проте головними
% викликами для хірургів стали:
% \begin{itemize}
%   \item довший і складніший процес навчання вишколювання хірургів;
%   \item вартість обладнання;
%   \item час операції довший за відкриту хірургію.
% \end{itemize}

% Якщо подивитися на весь процес то побачимо тенденцію у напрямку підвищення якості хірургії для пацієнта за рахунок
% ускладнення процесу для хірурга. Проте подальший розвиток технологій буде відбуватися у двох напрямках одночасно:
% якості
% для пацієнта і зручності для хірурга.

% Потенційні покращення, які варто запровадити:
% \begin{itemize}
%   \item здешевлення роботохірургії зокрема за рахунок пришвидшення процесу;
%   \item ефективніше вишколювання хірурга;
%   \item апробація нових інструментів для покращення часу і якості виконання хірургічної операції.
% \end{itemize}

% Важливим джерелом для дослідження вищезгаданих удосконалень є симуляційне середовище. У ньому хірург може
% вишколюватися.
% Крім того можна збирати інформацію про використання того чи іншого інструмента і апробовувати вдосконалення. Потенційно
% нові інструменти можна перевіряти у безпечному симульованому середовищі на багатьох сценаріях і порівнювати
% ефективність
% нового інструменту з уже наявними.

% Для більш реального моделювання процесів, які відбуваються у людському організмі, обирають методи, які враховують
% неоднорідність середовища та дозволяють обчислювати його напружено-деформований стан та температуру.

% Широко вживаними протягом кількох десятиріч методом розв'язування задач такого класу є проєкційно-сітковий
% метод (метод скінчених елементів, МСЕ). Завдяки таким властивостям як розрідженість і симетричність матриць для
% вказаних
% задач у локально-неоднорідних тілах різні варіанти МСЕ набули популярності.

% Із МСЕ цілком успішно може конкурувати, а іноді бути ефективнішим, велике сімейство методів граничних інтегральних
% рівнянь (МГІР) та граничних елементів (МГЕ), які базуються на ідеях розроблених С.Г.Міхліном, Н.І.Мусхелішвілі,
% В.Д.Купрадзе та іншими. Цій обставині особливо сприяють такі переваги, як зменшення на одиницю геометричної розмірності
% задачі, простота розв'язування для безмежних областей, можливість обчислювати у лінійних задачах значення шуканих
% величин у довільній точці без нової апроксимації, простота дискретизації.

% У свою чергу серед класу МГЕ можна виділити такі методи як комплексний метод граничних елементів (КМГЕ) і метод при
% граничних елементів (МПГЕ). На відміну від КМГЕ, який є видозміненим варіантом класичного МГЕ, МПГЕ базується на
% радикальніших ідеях. Порівняно із МГЕ, МПГЕ простіший для реалізації та у більшості випадків дає точніші результати при
% тих самих затратах обчислювальних ресурсів.

% Наукове завдання дисертаційної роботи -- розробка ефективного методу для розв'язування задач термопружності у
% композитах
% є актуальною проблемою.

% \textbf{Зв'язок роботи з науковими програмами, планами та темами.} Робота виконувалась на кафедрі програмування
% Львівського національного університету імені Івана Франка, а також у інституті прикладних проблем механіки і математики
% імені Ярослава Підстригача НАН України, відділ термомеханіки в рамках держбюджетної теми "Розробка аналітико-чисельних
% методів дослідження напруженого стану неоднорідних тіл з тепловими та залишковими деформаціями і дефектами структури"
% (номер державної реєстрації 0103U000131, 2003-2006 р.).

% \textbf{Мета роботи та задачі дослідження.} Метою роботи є: побудова нових підходів до розв'язування задач
% термопружності у композитах. Для досягнення поставленої мети необхідно вирішити такі завдання:
% \begin{itemize}
%   \item
%         дослідження ефективності нових методів
%   \item
%         порівняння швидкості збіжності і використання ресурсів із відомими методами
% \end{itemize}

% \emph{Об'єктом дослідження} є композити.

% \emph{Предметом дослідження} є нові методи розв'язування задач термопружності для композитів.

% \emph{Методи дослідження.} Для розв'язування поставлених задач використовується поєднання методів граничних і
% скінченних
% елементів.

% \textbf{Наукова новизна одержаних результатів.} Для тіл із локальними неоднорідностями побудований новий метод на
% основі
% МГЕ, МПГЕ і МСЕ для розв'язування задач:

% \begin{enumerate}
%   \item
%         теплопровідності
%   \item
%         пружності
%   \item
%         термопружності
% \end{enumerate}

% \textbf{Практичне значення отриманих результатів.} Отримані результати дозволяють ефективніше розв'язувати актуальні
% проблеми, які ставить перед нами економіка країни.

% \textbf{Апробація результатів дисертації.} Основні результати виконаних досліджень доповідалися та обговорювалися на:

% \begin{enumerate}
%   \def\labelenumi{\arabic{enumi}.}
%   \item
%         Восьмій Всеукраїнській науковій конференції (25-27 вересня 2001 р., м. Львів) ``Сучасні проблеми прикладної
%         математики
%         та інформатики'';
%   \item
%         Дев'ятій Всеукраїнській науковій конференції (24-26 вересня 2002 р., м. Львів) ``Сучасні проблеми прикладної
%         математики та інформатики'';
%   \item
%         Десятій Всеукраїнській науковій конференції (23-25 вересня 2003 р., м. Львів) ``Сучасні проблеми прикладної
%         математики
%         та інформатики'';
%   \item
%         Одинадцятій Всеукраїнській науковій конференції ({\hl !!! } 2004р., м. Львів) ``Сучасні проблеми прикладної
%         математики та інформатики'';
%   \item
%         Дванадцятій Всеукраїнській науковій конференції (4-6 жовтня 2005р., м. Львів) ``Сучасні проблеми прикладної
%         математики
%         та інформатики'';
%   \item
%         Науковому семінарі ІППММ НАНУ імені Я.С.Підстригача, присвяченої пам'яті Є.Г.Грицька (22 квітня 2003 р., м. Львів);
%   \item
%         Шостій Міжнародній науковій конференції (26-29 травня 2003 р., м. Львів) ``Математичні проблеми механіки неоднорідних
%         структур'';
%   \item
%         Семінар інституту прикладних проблем механіки та математики ім. Я. Підстригача
% \end{enumerate}

% У повному обсязі робота доповідалася на семінарі кафедри програмування Львівського національного університету імені
% Івана Франка.

% \textbf{Публікації.} Результати досліджень висвітлені в 4 статтях у журналах з переліку фахових видань ВАК України, в
% тому числі у доповідях НАН України.

% \textbf{Особистий внесок здобувача.} Всі результати дисертаційної роботи отримано здобувачем самостійно.

% \textbf{Структура та обсяг роботи.} Дисертація складається із вступу і шести розділів основної частини, висновків,
% списку використаної літератури.

% \begin{center}
%   \textbf{ОСНОВНИЙ ЗМІСТ РОБОТИ}
% \end{center}

% \textbf{Усталені домовленості.} У цій праці кома перед індексом означає похідну по змінній, яка відповідає цьому
% індексу, тобто ${f_{,j} = \partial}{f/\partial}x_{j}$, ${{f_{,ij} = \partial^{2}}{f/\partial}x_{i}\partial x_{j}}{}$,
% при цьому використовується німе сумування за індексами, які повторюються.

% У \textbf{вступі} наведено загальну характеристику роботи. Обґрунтовано актуальність теми, сформульовано мету та
% основні
% завдання дослідження. Визначено наукову новизну і практичну цінність отриманих результатів. Описано зв'язок роботи з
% науковими програмами та темами. Зазначено дані про особистий внесок автора, наведено дані про апробацію результатів
% роботи та публікації.

% У \textbf{першому розділі} міститься огляд праць стосовно сучасного стану досягнень у сфері математичного та
% комп'ютерного моделювання задач із локальними неоднорідностями. Також здійснено постановку задач дослідження. Основні
% ідеї математичного моделювання процесу термопружності базуються, зокрема, на роботах С.Г.Міхліна, Н.І.Мусхелішвілі,
% В.Д.Купрадзе, Є.Г. Грицька та інших.

% \textbf{Другий розділ} містить постановку стаціонарних задач теплопровідності, пружності і термопружності.

% \textbf{Стаціонарна задача теплопровідності.} Нехай нам задана область $\Omega$. Межу області позначимо $\Gamma$:
% $\Gamma = \partial \Omega$. Стаціонарне температурне поле у м'яких тканинах людського тіла є розв'язком рівняння
% Пеннеса:

% \begin{equation}
%   \label{eqn:theory_thermo_pennes_eqn}
%   \nabla \cdot \lambda_\theta \nabla \theta(x) + \omega \rho_b C_b (\theta_c - \theta(x)) + q_m  = 0, x \in \Omega
% \end{equation}

% \noindent де $\lambda_\theta$ -- коефіцієнт теплопровідності, $\theta$ -- шукана температура в області $\Omega$,
% $\rho_b$ -- густина крові, $\omega$ -- перфузія крові через орган, $C_b$ -- питома теплоємність крові,
% $\theta_c$ -- внутрішня температура тіла, $q_m$ -- метаболічне тепло у тканині. На межі задамо умови:

% \begin{equation}
%   \label{eqn:theory_thermo_cond_dirichlet}
%   \theta(x) = \theta_{\Gamma_D}(x), x \in \Gamma_D
% \end{equation}

% \begin{equation}
%   \label{eqn:theory_thermo_cond_neumann}
%   q(x) = q_{\Gamma_N}(x), x \in \Gamma_N
% \end{equation}

% \begin{equation}
%   \label{eqn:theory_thermo_cond_robin}
%   -\lambda_\theta q(x) = h(\theta(x) - \theta_e), x \in \Gamma_R,
% \end{equation}

% \noindent де $q(x) = \frac{\partial \theta(x)}{\partial n}$ -- тепловий потік, $n_i$ -- компоненти вектора зовнішньої
% нормалі до межі $\Gamma$, $\theta_e$ -- температура середовища, $h$ -- коефіцієнт тепловіддачі. Крайові умови
% (\ref{eqn:theory_thermo_cond_dirichlet}), (\ref{eqn:theory_thermo_cond_neumann}) і (\ref{eqn:theory_thermo_cond_robin})
% -- це умови Діріхле, Неймана і Робіна відповідно. Умова Робіна описує тепловий обмін тіла із навколишнім середовищем.

% \textbf{Стаціонарна задача пружності.} Розглянемо область $\Omega$. $u_i(x)$ -- позначимо компоненти вектора
% переміщень.
% Компоненти тензора деформацій $\varepsilon_{ij}$ визначаються співвідношеннями Коші:

% \begin{equation}
%   \label{eqn:theory_elasticity_gauchy_tensor}
%   \varepsilon_{ij}(x) = \frac{1}{2}(u_{i,j}(x) + u_{j,i}(x)), x \in \Omega
% \end{equation}

% Компоненти тензора напружень $\sigma_{ij}$ можна виразити через $\varepsilon_{ij}$ використавши закон Гука:

% \begin{equation}
%   \label{eqn:theory_elasticity_hooke_law}
%   \sigma_{ij}(x) = \lambda \delta_{ij}\varepsilon_{kk}(x) + 2\mu\varepsilon_{ij}(x), x \in \Omega,
% \end{equation}

% \noindent де $\lambda$ і $\mu$ -- це коефіцієнти Ляме, $\delta_{ij}$ -- символ Кронекера.

% Статична рівновага напружень і моментів виражається рівнянням:

% \begin{equation}
%   \label{eqn:theory_elasticity_sigma_equilibrium}
%   \sigma_{ij,i}(x) + b_j(x) = 0, x \in \Omega,
% \end{equation}

% \noindent де $b_j(x)$ -- компоненти вектора об'ємних напружень. Якщо у рівняння рівноваги
% (\ref{eqn:theory_elasticity_sigma_equilibrium}) підставити закон Гука (\ref{eqn:theory_elasticity_hooke_law}) і
% співвідношення Коші (\ref{eqn:theory_elasticity_gauchy_tensor}), то одержимо систему диференційних рівнянь у
% переміщеннях:

% \begin{equation}
%   \label{eqn:theory_elasticity_displacement_eqn}
%   \mu u_{i,jj}(x) + (\lambda + \mu)u_{j,ji} + b_i(x) = 0, x \in \Omega
% \end{equation}

% \noindent На межі задамо умови:
% \begin{equation}
%   \label{eqn:theory_elasticity_dirichlet}
%   u_i(x) = u_{\Gamma_Di}(x), x \in \Gamma_D
% \end{equation}

% \begin{equation}
%   \label{eqn:theory_elasticity_neumann}
%   t_i(x) = t_{\Gamma_Ni}(x), x \in \Gamma_N,
% \end{equation}

% \noindent де $t_i = \sigma_{ij}n_j$. Крайові умови (\ref{eqn:theory_elasticity_dirichlet}) і
% (\ref{eqn:theory_elasticity_neumann}) -- це умови Діріхле і Неймана відповідно.

% У практичних задачах замість коефіцієнтів Ляме часто задають модуль Юнга $E$ і коефіцієнт Пуассона $\nu$. Для
% перетворення цих величин використовують формули:

% \begin{equation}
%   \label{eqn:theory_elasticity_lame_lambda_convertion}
%   \lambda = \frac{E\nu}{(1+\nu)(1-2\nu)}
% \end{equation}

% \begin{equation}
%   \label{eqn:theory_elasticity_lame_mu_convertion}
%   \mu = \frac{E}{2(1+\nu)}
% \end{equation}

% \textbf{Стаціонарна задача термопружності.} Розглянемо вплив температури на пружність м'яких тканин тіла.
% Представимо тіло областю $\Omega$. До закону Гука (\ref{eqn:theory_elasticity_hooke_law}) додамо співвідношення
% Дюгамеля-Неймана і отримаємо:

% \begin{equation}
%   \label{eqn:theory_thermoelasticity_duhamel_neumann_extension}
%   \sigma_{ij}(x) = \lambda \delta_{ij}\varepsilon_{kk}(x) + 2\mu\varepsilon_{ij}(x) - \beta \theta, x \in \Omega,
% \end{equation}

% \noindent де $\beta = (2\mu + 3\lambda)\alpha$, $\alpha$ -- коефіцієнт теплового розширення матеріалу. Температурне
% поле
% $\theta$ знайдемо із рівняння (\ref{eqn:theory_thermo_pennes_eqn}). Запишемо рівняння рівноваги у переміщеннях із
% використанням (\ref{eqn:theory_thermoelasticity_duhamel_neumann_extension}):

% \begin{equation}
%   \label{eqn:theory_thermoelasticity_displacement_eqn}
%   \mu u_{i,jj}(x) + (\lambda + \mu)u_{j,ji} - \beta \theta_{,i} + b_i(x) = 0, x \in \Omega
% \end{equation}

% На межі $\Gamma$ можна задати умови для рівняння Пеннеса (\ref{eqn:theory_thermo_cond_dirichlet}),
% (\ref{eqn:theory_thermo_cond_neumann}) і (\ref{eqn:theory_thermo_cond_robin}), а для переміщень і поверхневих напружень

% (\ref{eqn:theory_elasticity_dirichlet}) і (\ref{eqn:theory_elasticity_neumann}) відповідно.

% У \textbf{третьому розділі} сформульовано плоску і просторову задачі стаціонарної теплопровідності із злоякісними
% включеннями (\ref{eqn:theory_thermo_pennes_eqn}), (\ref{eqn:theory_thermo_cond_dirichlet}),
% (\ref{eqn:theory_thermo_cond_neumann}) і (\ref{eqn:theory_thermo_cond_robin}). Коефіцієнт теплопровідності залежить від

% координат в області вишколювання і сталий поза нею:

% \begin{equation}
%   \label{eqn:thermo_thermal_conductivity}
%   \lambda_\theta(x) = \lambda_{\theta H}(x)\chi_H + \lambda_{\theta c}, x \in \Omega
% \end{equation}

% \begin{equation*}
%   \lambda_{\theta H}(x) = (0,495 - \lambda_{\theta c}) \cdot cos\left[\frac{\pi}{2}\frac{||x-x_H||^2}{r_H^2}\right]
% \end{equation*}

% \begin{equation*}
%   \lambda_{\theta c} = 0,19;
% \end{equation*}

% \noindent де $\chi_H$ - характеристична функція області $H$, a $x_H$ і $r_H$ --- центр області і радіус відповідно.
% Для подальших міркувань буде важливим, що $\lambda_{\theta H}(x) = 0$ для всіх $x \in L$, де $L = \partial H$.

% Задача розв'язується методом адитивного розщеплення оператора рівняння (\ref{eqn:theory_thermo_pennes_eqn}) на дві
% частини: зі сталим коефіцієнтом теплопровідності $\lambda_{\theta c}$ (\ref{eqn:thermo_thermal_conductivity}) і з
% залежним від координат $\lambda_{\theta H}(x)$ лише в області $H$:

% \begin{equation}
%   \label{eqn:thermo_pennes_separation_eqn}
%   \lambda_{\theta c} \nabla^2 \theta(x) + \omega \rho_b C_b (\theta_c - \theta(x)) =
%   - \nabla \lambda_{\theta H} (x) \cdot \nabla \theta(x) \cdot \chi_H, x \in \Omega
% \end{equation}

% Для МГЕ інтегральне подання розв'язку рівняння (\ref{eqn:thermo_pennes_separation_eqn}) запишемо у вигляді:

% \begin{equation}
%   \label{eqn:thermo_bem_direct_pennes}
%   \begin{split}
%     c(\xi)\theta(\xi) + \int \limits_{\Gamma}\theta(x)\hat{q}(x,\xi)d\Gamma(x)
%     + \int \limits_{H} \nabla{\lambda_{\theta H}(x)}\nabla{\theta(x)} \hat{\theta}(x,\xi) d H(x) = \\
%     = \int \limits_{\Gamma}q(x)\hat{\theta}(x,\xi)d\Gamma(x),
%   \end{split}
% \end{equation}

% \noindent де $\hat{\theta}$ і $\hat{q}$ -- фундаментальні розв'язки рівняння (\ref{eqn:theory_thermo_pennes_eqn}). За
% допомогою теореми Гріна рівняння (\ref{eqn:thermo_bem_direct_pennes}) зводимо до:

% \begin{equation}
%   \label{eqn:thermo_bem_direct_pennes_simplified}
%   \begin{split}
%     c(\xi)\theta(\xi) + \int \limits_{\Gamma}\theta(x)\hat{q}(x,\xi)d\Gamma(x) = \\
%     = \int \limits_{H}\theta(x) [ \nabla^2 \lambda_{\theta H}(x) \hat{\theta}(x,\xi) +
%       \nabla \lambda_{\theta H}(x) \nabla \hat{\theta}(x,\xi) ] dH(x)
%     + \int \limits_{\Gamma}q(x)\hat{\theta}(x,\xi)d\Gamma(x)
%   \end{split}
% \end{equation}

% Для МПГЕ розглянемо таку область $B \subset R^n$, що $\Omega \subset B$, $\partial\Omega \cap \partial B =
%   \varnothing$.
% Позначимо приграничну область $G = B \setminus \Omega$. Потім застосуємо теорему Гріна до інтегрального подання і
% отримаємо:

% \begin{equation}
%   \label{eqn:thermo_cobem_pennes}
%   \int \limits_G g(x)\hat{\theta}(x,\xi)dG(x)
%   = \int \limits_{H}\theta(x) [ \nabla^2 \lambda_{\theta H}(x) \hat{\theta}(x,\xi) +
%     \nabla \lambda_{\theta H}(x) \nabla \hat{\theta}(x,\xi) ] dH(x)
% \end{equation}

% Із інтегральних подань розв'язків (\ref{eqn:thermo_bem_direct_pennes_simplified}) для МГЕ і
% (\ref{eqn:thermo_cobem_pennes}) для МПГЕ отримаємо систему лінійних алгебраїчних рівнянь для значень температури і
% потоку на межі, а також для температури всередині області $H$. Цю область розіб'ємо на скінченні елементи, а $G$ на
% приграничні у випадку МПГЕ.

% Завдяки такому підходу отримано розв'язки для включень різних розмірів і відстаней від поверхні органу. Вплив
% розташування і радіуса злоякісного утворення на температурне поле здорової тканини узгоджується із клінічними
% спостереженнями. Цей підхід можна використовувати для моделювання температури злоякісних включень у медичних
% симуляторах.

% \textbf{Четвертий розділ} присвячений задачі пружності м'яких тканини людини із неоднорідностями. Обчислюється тиск,
% який чинить тканина на медичний інструмент під час контакту.

% \begin{figure}[ht!]
%   \centering
%   \begin{tikzpicture}
%     \draw (5.7,3) rectangle (6.3,7);
%     \draw (1,1) rectangle (11,3);
%     \draw (6,2.5) circle (0.5);
%     \draw[thick, ->] (5.3,6) -- (5.3,4);

%     \draw[thick,->] (0,0) -- (12,0) node[anchor=north west] {x};
%     \draw[thick,->] (0,0) -- (0,9) node[anchor=south east] {y};

%     \foreach \x in {0,...,4}
%     \pgfmathtruncatemacro{\label}{\x * 10}
%     \draw ( \x * 2.5 cm,1pt) -- ( \x * 2.5 cm,-1pt) node[anchor=north] {\label};
%     \foreach \y in {0,...,3}
%     \pgfmathtruncatemacro{\label}{\y * 10}
%     \draw (1pt,\y * 2.5 cm) -- (-1pt,\y * 2.5 cm) node[anchor=east] {\label};

%     \node at (6,3.5) {$\Gamma_i$};

%     \node at (8,2) {$\Omega$};
%     \node at (6,5) {$\Omega_i$};
%     \node at (6,2.5) {$H$};
%     \node at (4,0.5) {$\Gamma_b$};

%   \end{tikzpicture}

%   \caption{Переріз тканин і інструмента}
%   \label{fig:elasticity_domain}
% \end{figure}

% Як зображено на рис. \ref{fig:elasticity_domain} на частині межі $\Gamma_i$ медичний інструмент взаємодіє із
% органом людини $\Omega$, що містить включення $H$. Для знаходження тиску органу на інструмент розв'яжемо задачу
% (\ref{eqn:theory_elasticity_displacement_eqn}), (\ref{eqn:theory_elasticity_dirichlet}) і
% (\ref{eqn:theory_elasticity_neumann}). Коефіцієнти Ляме сталі на всій області $\Omega$ крім $H$, де вони залежать
% від координат:

% \begin{equation}
%   \label{eqn:elasticity_lyame_lambda}
%   \lambda(x) = \lambda_H(x)\chi_H + \lambda_c, x \in \Omega
% \end{equation}

% \begin{equation*}
%   \lambda_H(x) = (3714,3 - \lambda_c) \cdot cos\left[\frac{\pi}{2}\frac{||x-x_H||^2}{r_H^2}\right]
% \end{equation*}

% \begin{equation*}
%   \lambda_c = 660,1
% \end{equation*}

% \begin{equation}
%   \label{eqn:elasticity_lyame_mu}
%   \mu(x) = \mu_H(x)\chi_H + \mu_c, x \in \Omega
% \end{equation}

% \begin{equation*}
%   \mu_H(x) = (3428,6 - \mu_c)cos\left[\frac{\pi}{2}\frac{||x-x_H||^2}{r_H^2}\right]
% \end{equation*}

% \begin{equation*}
%   \mu_c = 404,6;
% \end{equation*}

% \noindent де $\chi_H$ --- характеристична функція області $H$, a $x_H$ і $r_H$ --- центр і радіус області відповідно.
% Надалі буде важливим, що $\lambda_H(x) = 0$ і $\mu_H(x) = 0$ для всіх $x \in L$, де $L = \partial H$.

% Застосуємо адитивне розщеплення оператора рівняння (\ref{eqn:theory_elasticity_displacement_eqn}) на дві частини: із
% сталими $\lambda_c$ і $\mu_c$ (\ref{eqn:elasticity_lyame_lambda}), (\ref{eqn:elasticity_lyame_mu}) і залежними від
% координат $\lambda_H(x)$ і $\mu_H(x)$ -- лише в області $H$:

% \begin{equation}
%   \label{eqn:elasticity_displacement_separation_eqn}
%   \mu_c u_{i,jj}(x) + (\lambda_c + \mu_c)u_{j,ji} = -f_{ij,j}(\varepsilon(x)), x \in \Omega,
% \end{equation}

% \noindent де

% \begin{equation*}
%   f_{ij,j}(\varepsilon(x)) = [ 2\mu_H(x)\varepsilon_{ij}(x) + \lambda_H(x) \delta_{ij} \varepsilon_{kk}(x) ] \chi_H
% \end{equation*}

% Для МГЕ інтегральне подання розв'язку рівняння (\ref{eqn:elasticity_displacement_separation_eqn}) запишемо у вигляді:

% \begin{equation}
%   \label{eqn:elasticity_bem_direct_elasticity}
%   \begin{split}
%     c_{ij}(\xi) u_j(\xi) = \int \limits_{\Gamma}\hat{u}_{ij}(x,\xi)t_j(x)d\Gamma(x) -
%     \int \limits_{\Gamma}\hat{t}_{ij}(x,\xi)u_j(x)d\Gamma(x) + \\
%     + \int \limits_H \hat{u}_{ij}(x,\xi)f_{jk,k}(\varepsilon(x))dH(x)
%   \end{split}
% \end{equation}

% \noindent$\hat{u}_{ij}$ і $\hat{t}_{ij}$ -- фундаментальні розв'язки рівняння пружності
% (\ref{eqn:theory_elasticity_displacement_eqn}) для переміщень і напружень відповідно. Застосуємо до
% (\ref{eqn:elasticity_bem_direct_elasticity}) теорему Гріна для інтегрування частинами. Одержимо:

% \begin{equation}
%   \label{eqn:elasticity_bem_direct_elasticity_simplified}
%   \begin{split}
%     c_{ij}(\xi) u_j(\xi) = \int \limits_{\Gamma}\hat{u}_{ij}(x,\xi)t_j(x)d\Gamma(x) -
%     \int \limits_{\Gamma}\hat{t}_{ij}(x,\xi)u_j(x)d\Gamma(x) - \\
%     - \int \limits_H \hat{u}_{ij,k}(x,\xi)f_{jk}(\varepsilon(x))d H(x)
%   \end{split}
% \end{equation}

% У випадку МПГЕ із приграничною областю $G$ застосуємо теорему Гріна до інтегрального подання і отримаємо:

% \begin{equation}
%   \label{eqn:elasticity_cobem_direct_elasticity_simplified}
%   \begin{split}
%     u_i(\xi) = \int \limits_G \hat{u}_{ij}(x,\xi) g_j(x)d G(x)
%     - \int \limits_H \hat{u}_{ij,k}(x,\xi)f_{jk}(\varepsilon(x))d H(x)
%   \end{split}
% \end{equation}

% Із інтегральних подань розв'язків (\ref{eqn:elasticity_bem_direct_elasticity_simplified}) для МГЕ і
% (\ref{eqn:elasticity_cobem_direct_elasticity_simplified}) для МПГЕ отримаємо систему лінійних алгебраїчних рівнянь для
% значень переміщень і поверхневих напружень на межі, а також для елементів тензора деформацій всередині області $H$. Цю
% область розіб'ємо на скінчені елементи, а $G$ на приграничні у випадку МПГЕ.

% Завдяки такому підходу отримано розв'язки для включень різних розмірів і відстаней від поверхні органу. Тиск із
% яким м'яка тканина із включенням діє на інструмент узгоджується із клінічними спостереженнями. Цей підхід можна
% використовувати для розрахунку тиску у медичних симуляторах.

% У \textbf{п'ятому розділі} розглядається задачі термопружності м'яких тканини людини із неоднорідностями. Обчислюється
% тиск, який чинить тканина на медичний інструмент під час контакту із урахуванням теплових процесів.

% Як і для задачі пружності умови задачі відображені на рис. \ref{fig:elasticity_domain}. Інструмент $\Omega_i$ на
% частині
% межі $\Gamma_i$ взаємодіє із органом людини $\Omega$, що містить включення $H$. Теплові процеси в органі описуються
% рівнянням Пеннеса (\ref{eqn:theory_thermo_pennes_eqn}).

% Для знаходження тиску органу на інструмент розв'яжемо систему рівнянь (\ref{eqn:theory_thermo_pennes_eqn}) і
% (\ref{eqn:theory_thermoelasticity_displacement_eqn}) з умовами (\ref{eqn:theory_thermo_cond_dirichlet}),
% (\ref{eqn:theory_thermo_cond_neumann}) і (\ref{eqn:theory_thermo_cond_robin}) та
% (\ref{eqn:theory_elasticity_dirichlet})
% і (\ref{eqn:theory_elasticity_neumann}). Коефіцієнти теплопровідності і Ляме сталі у всій області $\Omega$ крім $H$, де
% вони залежать від координат за формулами (\ref{eqn:thermo_thermal_conductivity}), (\ref{eqn:elasticity_lyame_lambda}) і
% (\ref{eqn:elasticity_lyame_mu}).

% Застосуємо адитивне розщеплення операторів рівнянь (\ref{eqn:theory_thermo_pennes_eqn}) і
% (\ref{eqn:theory_thermoelasticity_displacement_eqn}) на частини із сталими $\lambda_{\theta c}$, $\lambda_c$ і $\mu_c$
% (\ref{eqn:thermo_thermal_conductivity}), (\ref{eqn:elasticity_lyame_lambda}), (\ref{eqn:elasticity_lyame_mu}) і
% залежними від координат $\lambda_{\theta H}(x)$, $\lambda_H(x)$ і $\mu_H(x)$ лише в області $H$. Отримаємо рівняння
% (\ref{eqn:thermo_pennes_separation_eqn}) і рівняння для переміщень:

% \begin{equation}
%   \label{eqn:thermoelasticity_displacement_separation_eqn}
%   \mu_c u_{i,jj}(x) + (\lambda_c + \mu_c)u_{j,ji} = -f_{ij,j}(\varepsilon(x)) + \beta_{,j}(x) \theta(x) +
%   \beta(x) \theta_{,j}(x) , x \in \Omega,
% \end{equation}

% \noindent де

% \begin{equation*}
%   f_{ij,j}(\varepsilon(x)) = [ 2\mu_H(x)\varepsilon_{ij}(x) + \lambda_H(x) \delta_{ij} \varepsilon_{kk}(x) ] \chi_H
% \end{equation*}

% Для МГЕ інтегральне подання розв'язку системи рівнянь (\ref{eqn:thermo_pennes_separation_eqn}),
% (\ref{eqn:thermoelasticity_displacement_separation_eqn}) запишемо як (\ref{eqn:thermo_bem_direct_pennes}) і додамо:

% \begin{equation}
%   \label{eqn:elasticity_bem_direct_thermoelasticity}
%   \begin{split}
%     c_{ij}(\xi) u_j(\xi) = \int \limits_{\Gamma}\hat{u}_{ij}(x,\xi)t_j(x)d\Gamma(x) -
%     \int \limits_{\Gamma}\hat{t}_{ij}(x,\xi)u_j(x)d\Gamma(x) + \\
%     + \int \limits_H \hat{u}_{ij}(x,\xi)f_{jk,k}(\varepsilon(x))dH(x)
%     - \int \limits_{\Omega} \hat{u}_{ij}(x,\xi)[\beta_{,j}(x)\theta(x)+\beta(x)\theta_{,j}(x)]d\Omega(x)
%   \end{split}
% \end{equation}

% \noindent$\hat{u}_{ij}$ і $\hat{t}_{ij}$ -- фундаментальні розв'язки рівняння пружності
% (\ref{eqn:theory_elasticity_displacement_eqn}) для переміщень і напружень відповідно. За допомогою теореми Гріна
% систему
% рівнянь (\ref{eqn:thermo_bem_direct_pennes})(\ref{eqn:elasticity_bem_direct_thermoelasticity}) можна звести до
% (\ref{eqn:thermo_bem_direct_pennes_simplified}) і:

% \begin{equation}
%   \label{eqn:elasticity_bem_direct_thermoelasticity_simplified}
%   \begin{split}
%     c_{ij}(\xi) u_j(\xi) = \int \limits_{\Gamma}\hat{u}_{ij}(x,\xi)t_j(x)d\Gamma(x) -
%     \int \limits_{\Gamma}\hat{t}_{ij}(x,\xi)u_j(x)d\Gamma(x) - \\
%     - \int \limits_H \hat{u}_{ij,k}(x,\xi)f_{jk}(\varepsilon(x))d H(x)
%     - \int \limits_{\Gamma} \hat{u}_{ij} \beta(x) \theta (x) n_j d \Gamma(x) + \\
%     + \int \limits_{\Omega} \theta(x) [\hat{u}_{ij,j}(x,\xi)\beta(x)-\hat{u}_{ij}(x,\xi)\beta_{,j}(x)]d\Omega(x)
%   \end{split}
% \end{equation}

% Для МПГЕ після застосування теореми Гріна до інтегрального подання системи рівнянь отримаємо
% (\ref{eqn:thermo_cobem_pennes}) і рівняння:

% \begin{equation}
%   \label{eqn:thermoelasticity_cobem_direct_elasticity_simplified}
%   \begin{split}
%     u_i(\xi) = \int \limits_G \hat{u}_{ij}(x,\xi) g_j(x)d G(x)
%     - \int \limits_H \hat{u}_{ij,k}(x,\xi)f_{jk}(\varepsilon(x))d H(x) - \\
%     - \int \limits_{\Gamma} \hat{u}_{ij} \beta(x) \theta (x) n_j d \Gamma(x)
%     + \int \limits_{\Omega} \theta(x) [\hat{u}_{ij,j}(x,\xi)\beta(x)-\hat{u}_{ij}(x,\xi)\beta_{,j}(x)]d\Omega(x)
%   \end{split}
% \end{equation}

% Із інтегральних подань розв'язків системи рівнянь
% (\ref{eqn:thermo_bem_direct_pennes_simplified})(\ref{eqn:elasticity_bem_direct_elasticity_simplified}) для МГЕ і
% (\ref{eqn:thermo_cobem_pennes})(\ref{eqn:elasticity_cobem_direct_elasticity_simplified}) для МПГЕ отримаємо систему
% лінійних алгебраїчних рівнянь для значень температурного поля, переміщень і поверхневих напружень на межі, а також для
% елементів тензора деформацій всередині області $H$. Цю область розіб'ємо на скінчені елементи, а $G$ на приграничні у
% випадку МПГЕ.

% Отримані розв'язки задачі термопружності дозволяють оцінити температурне поле і поверхневі напруження у медичному
% симуляторі. Отримані дані узгоджуються із результатами спостережень описаних у медичних наукових працях.

% \textbf{Шостий розділ} містить інформацію про програмне середовище розроблене для проведення наукових досліджень, а
% також для викладення отриманих результатів у зручній формі. Технічне рішення дозволяє створювати відтворювані
% налаштування робочого оточення і контрольовано вносити до нього зміни.Розроблена система може легко розповсюджуватися
% між різними машинами, що забезпечує високу гнучкість у роботі. Модульність підходу у рішенні полегшує локалізацію
% некоректної поведінки системи і незалежне відлагодження її частин. Більшість бібліотек і утиліт, яке використовуються у
% програмному середовищі знаходяться у вільному доступі у мережі Інтернет.

% \begin{center}
%   \textbf{ОСНОВНІ РЕЗУЛЬТАТИ ТА ВИСНОВКИ}
% \end{center}

% Дисертаційна робота є завершеним науковим дослідженням, у якому розв'язано важливе науково-технічне завдання
% математичного моделювання процесів термопружності у м'яких тканинах тіла людини. У моделях фізичні характеристики у
% областях неоднорідностей можуть залежати від координат.

% При цьому отримано наступні наукові та практичні результати:

% \begin{enumerate}
%   \item Побудовано математичну модель теплообміну в органах людського тіла із урахуванням злоякісних включень під час
%         процедури малоінвазивної роботохірургії. Параметри задачі отримано із міжнародних медичних баз даних. Побудовані
%         числові розв'язки відповідних крайових задач, обґрунтованість яких забезпечується високим рівнем узгодження
%         результатів числових експериментів розв'язками відповідних типів модельних тестових задач. На основі отриманих
%         числових результатів підтверджено медичні спостереження про процеси теплообміну у органах людини із злоякісними
%         включеннями.
%   \item Вперше використано МГЕ і МПГЕ у поєднанні із скінченними елементами для розв'язування задачі теплообміну у
%         людському тілі із локальними неоднорідностями.
%   \item Досліджено вплив злоякісного включення в органі на поверхневі напруження органу. Обчислена сила із якою орган
%         діє на інструмент робота хірурга, який перебуває у жорсткому контакті із його поверхнею. Отримані числові результати
%         узгоджуються із результатами із клінічних досліджень описаних у наукових працях.
%   \item Використано МГЕ і МПГЕ у поєднанні із скінченними елементами для розв'язування задачі пружності у людському
%         тілі
%         із локальними неоднорідностями. Злоякісні утворення моделювалися коефіцієнтами Ляме, які були сталими скрізь, крім
%         області неоднорідності. Такий підхід дозволи обмежитися скінченними елементами тільки в області злоякісного
%         включення.
%   \item Побудована модель теплового процесу в органі людського тіла із пружною взаємодією із інструментом робота
%         хірурга. Чисельно розв'язана відповідна крайова система рівнянь. Отримані результати допоможуть створювати ефективні
%         симуляції для хірургів, а також досліджувати нові інструменти і зручність їх використання під час операцій.
% \end{enumerate}

% \begin{center}
%   \textbf{СПИСОК ОПУБЛІКОВАНИХ ПРАЦЬ ЗА ТЕМОЮ ДИСЕРТАЦІЇ}
% \end{center}

% Гудзь Р., Петльований А. Моделювання стаціонарного процесу теплопровідності у тілі зі сукупністю локальних
% неоднорідностей, застосовуючи комплексний метод граничних елементів // Вісн. Львів. Ун-ту. Сер. прикл. матем. та
% інформ.
% –2003. – Вип. 6. – С. 144 – 150

% Гудзь Р., Петльований А. Комплексний метод граничних елементів при моделюванні фізичних процесів у тілах з
% композиційних
% матеріалів // Вісн. Львів. Ун-ту. Сер. прикл. матем. та інформ. –2003. – Вип. 7. – С. 148 – 155

% Гудзь Р., Журавчак Л., Петльований А. Моделювання напружено-деформованого  стану кусково-однорідних тіл методами
% граничних та приграничних елементів // Вісн. Львів. Ун-ту. Сер. прикл. матем. та інформ. - 2004. –Вип 9, – C. 145-156

% Гудзь Р. В., Журавчак Л. М., Петльований А. Т. Поєднання граничних, приграничних та скінченних елементів в задачі
% пружності для неоднорідного тіла // Доповіді НАН України. – 2005. – №. 5.– С. 46-50.

% Гудзь Р. В., Журавчак Л. М., Петльований А. Т. Розв’язування статичної задачі термо-пружності для
% локально-неоднорідного
% тіла поєднанням граничних, приграничних та скінченних елементів // Мат. методи та фізико-механічні поля, 2006.–49, №
% 2.–
% С. 148-156.

% Гудзь Р. В., Журавчак Л. М., Петльований А. Т. Порівняння методів граничних та пригранич­них елементів при моделюванні
% напружено-деформованого стану кусково-однорідних об’єктів // Тези доповідей X Всеукраїнської. наук. конф. “Сучасні
% проблеми прикладної математики та інформатики”(23–25 вересня 2003 р., м. Львів). – Львів, Видавн. центр. ЛНУ ім. І.
% Франка, 2003. – С. 47.

% Гудзь Р. В., Журавчак Л. М., Петльований А. Т. Плоска задача термо­пружності для компози­ту // Тези доп. Всеукр. наук.
% конф. “Сучасні проблеми прикладної математики та інформатики” (21-23 вересня 2004 р.).– Львів, 2004. – С.59–60.

% Гудзь Р. В., Журавчак Л. М., Петльований А. Т. Поєднання граничних, приграничних та скін­ченних елемен­тів  у задачі
% термопружності для локально-неоднорідного тіла // Тез. доп. конф. молодих учених із су­часних проблем механіки і
% математики імені академіка Я.С. Підстригача.– Львів, 2005.– C. –.

% \begin{center}
%   \textbf{АНОТАЦІЇ}
% \end{center}

% \textbf{Петльований А.Т. Моделювання процесів термопружності у композитах методами граничних і приграничних елементів.}
% -- На правах рукопису.

% Дисертація на здобуття наукового ступеня кандидата технічних наук за спеціальністю 01.05.02 --
% математичне моделювання та обчислювальні методи. Львівський національний університет імені Івана Франка, Міністерство
% освіти і науки України, Львів, 2025.

% Дисертація присвячена математичному моделюванню процесів термопружності у композитах з урахуванням локальних
% неоднорідностей, які залежать від координат. Для цього використовується поєднання методів граничних і приграничних
% елементів із методом скінчених елементів, але тільки в області неоднорідності.

% \textit{Ключові слова}: МГЕ, МПГЕ, термопружність, рівняння Пеннеса, м'які тканини тіла, локальні неоднорідності,
% математичне моделювання.

% \textbf{Petlovanyy A.T. Using boundary and co-boundary elements methods to model thermoelasticity processes in
%   composite
%   materials.} -- Manuscript.

% Thesis for PhD degree on technical sciences in specialty 01.05.02 -- mathematical modeling and computational methods.
% Ivan Franko National University of Lviv, Ministry of Education and Science of Ukraine, Lviv, 2025.

% The thesis is devoted to mathematical modeling of processes of thermoelasticity in composite materials with spatial
% inhomogeneity. BEM and CBEM are used together with FEM (only in are of inhomogeneity) to solve these types of problems.

% \textit{Keywords}: BEM, CBEM, thermoelasticity, Pennes equation, human tissue, inhomogeneous inclusion,
% mathematical modeling.
